\section{Introducción}
\label{sec:introduccion}

Un mapa de cultivos por parcela promete un catálogo. El catálogo es una decisión: alguien eligió
qué tipos de cultivo va a nombrar el producto y cuáles no, y esa elección se toma mucho antes de
que se calcule ninguna cifra de calidad. En la práctica rara vez se reporta como una elección.
Llega en una frase sobre que algunas clases tenían muy poca muestra, y la evaluación que sigue se
calcula sobre lo que quedó.

Esto importa porque la métrica que domina la literatura de mapeo de cultivos desbalanceado es el
F1 promediado por clase, y promediar por clase divide entre el número de clases. Retirar una clase
que el modelo manejaba mal sube el promedio dos veces: una al borrar del numerador un F1 bajo, y
otra al encoger el denominador. El segundo efecto no requiere modelo, ni mecanismo, ni trabajo
alguno. Es aritmética. Y sin embargo, una mejora producida enteramente por el segundo efecto es
indistinguible, en una tabla de resultados, de una mejora producida por un método.

Tomamos esa ambigüedad como objeto de estudio. Consideramos dos mecanismos con los que se puede
cambiar cobertura por calidad macro. El primero recorta la leyenda prometida: el producto nombra
$K$ de las $C$ clases disponibles y declina responder allí donde su predicción sin restringir cae
fuera de ellas. El segundo conserva la leyenda completa y se abstiene en las parcelas donde el
predictor está menos seguro, en la tradición de la clasificación selectiva
\cite{chow1970reject,geifman2017selective}. Los dos reducen cuánto mapa se entrega, los dos suben
la cifra macro, y la pregunta es de cuánto de esa subida es responsable cada uno.

Nuestra respuesta es un control, y es casi vergonzosamente barato. Puntúese al predictor
\emph{intacto} sobre esa misma leyenda recortada, entregándolo todo. Lo que ese control recupere
nunca lo compró el mecanismo. En un primer banco público de dieciocho clases agronómicas, el
control se lleva el \num{88.3}\,\% de la mejora aparente al pasar de dieciocho clases a ocho. En
un segundo banco, de otro año, otra región y otra familia de modelo, se lleva el \num{95.1}\,\% en
el punto de operación preregistrado. Y en el tramo donde la cobertura sigue completa se lleva
\emph{toda}: cada mecanismo que probamos devuelve la misma cifra, hasta el cuarto decimal, mientras
el promedio macro sube \num{0.2447}.

Del mismo protocolo salen otros dos resultados. A igual cobertura, abstenerse en las parcelas de
baja confianza supera al recorte de leyenda en los dos bancos, que es lo contrario de lo que
supone la práctica que nos propusimos formalizar. Y retirar clases por soporte bajo, que es el
criterio que se reporta usar, es el más débil de los tres que probamos y no es monótono en el
número de clases retiradas: retirar más puede empeorar el mapa entregado.

Las contribuciones son las siguientes. Formalizamos los dos mecanismos y damos un protocolo de
bloques espaciales que elige el punto de operación fuera de los datos que lo miden, puntúa ambos
mecanismos sobre el mismo conjunto de clases para que sus cifras sean comparables, y nunca consulta
una etiqueta verdadera para decidir qué parcelas entrega. Cuantificamos la descomposición de la
mejora aparente en denominador y mecanismo sobre dos bancos públicos bajo un protocolo que no
cambia entre uno y otro. Ordenamos los tres criterios de retirada con intervalos pareados y
corrección por multiplicidad. Y mostramos que un meta-modelo de apilamiento entrenado ejecuta esa
misma retirada de forma implícita y no declarada, y que la diferencia entre medirla dentro y fuera
de muestra basta para invertir la conclusión sobre si el ensamble aporta algo.
